\begin{recipe}{Gnocchi met salie-roomboter}
\kicker[Hoofdgerecht,Vegetarisch,Italiaans]{Hoofdgerecht · vegetarisch · Italiaans}
\index[register]{Gnocchi met salie-roomboter}
\index[register]{Gnocchi}
\index[register]{Salie}
\index[register]{Parmezaanse kaas}

\meta{15 minuten}

\lettrine{V}{erse} gnocchi met gebruinde boter en salie, in Italië bekend als
burro e salvia, is zo klaar, en een eenvoudige combinatie die goed werkt. De
boter geeft een nootachtige smaak, de salie wordt krokant.

\blockrule{Ingrediënten}
\begin{ingredients}
  \ing{2 pakken}{verse gnocchi}\index[register]{Gnocchi}
  \ing{50 g}{roomboter}
  \ing{10 blaadjes}{salie}\index[register]{Salie}
  \ingb{een paar druppels citroensap}
  \ingb{Parmezaanse kaas, geraspt}\index[register]{Parmezaanse kaas}
  \ingb{versgemalen peper en zout}
\end{ingredients}

\blockrule{Bereiding}
\tip{Let op bij het bruinen van de boter: zodra hij nootachtig ruikt en de
salie krokant is, moet hij van het vuur. Daarna verbrandt hij snel.}
\begin{steps}
  \item Kook de gnocchi volgens de aanwijzingen op de verpakking (of je
        eigen gnocchi, p.~\pageref{recipe:gnocchi}) in ruim gezouten water.
        Schep vlak voor het afgieten een lepel kookvocht apart.
  \item Smelt ondertussen de roomboter in een hapjespan op middelhoog vuur
        en voeg de salieblaadjes toe. Laat de boter goudbruin en nootachtig
        worden en roer er een paar druppels citroensap door: dat geeft het
        nootachtige vet een frisse rand.
  \item Zodra de gnocchi bovendrijven, schep je ze met een schuimspaan
        rechtstreeks in de pan met de salieboter. Voeg de lepel kookvocht toe
        en beweeg de pan: boter en zetmeelwater emulgeren tot een glanzend
        sausje in plaats van een laagje vet.
  \item Schep om tot de gnocchi bedekt zijn met de saus en breng op smaak
        met peper en zout.
  \item Serveer direct met een flinke schep Parmezaanse kaas erover.
\end{steps}
\end{recipe}
